% Môn: Vật lý
% Lớp: 12
% Chương: Chương I. Vật lí nhiệt
% Bài: L12C1 Bài 1. Cấu trúc của chất. Sự chuyển thể · Quá trình chuyển pha
% ID: L12C1B3-121-TN
% ID: L12C1B3-121-TN
% Mức: NB
% ID: L12C1B3-121-TN


% Mức: TH
% Mức: TH
% ID: L12C1B3-124-TN
% ID: L12C1B3-124-TN
% Mức: TH
% ID: L12C1B3-124-TN
% ID: L12C1B3-124-TN
% Mức: TH
% ID: L12C1B1-01-TN
%=== Câu 4 ===%
\dangbt{Quá trình chuyển pha}
\begin{ex}
% ID: L12C1B1-01-TN
% Mức: TH
	Trong quá trình hóa hơi của một lượng chất lỏng ở nhiệt độ sôi thì
	\choice 
	{\True nhiệt độ chất lỏng không thay đổi}
	{thể tích khối chất lỏng không thay đổi}
	{nhiệt độ của vật tăng liên tục}
	{nhiệt độ của chất lỏng giảm liên tục}
	\loigiai{
		Khi chất lỏng hóa hơi ở đúng nhiệt độ sôi, nhiệt lượng cung cấp cho chất lỏng được dùng để thực hiện sự chuyển thể từ lỏng sang hơi, chứ không làm tăng nhiệt độ.
		
		Vì vậy trong suốt quá trình hóa hơi ở nhiệt độ sôi, \textbf{nhiệt độ của chất lỏng không thay đổi}.
	}
\end{ex}

% Mức: TH
% Mức: TH
% ID: L12C1B3-125-TN
% ID: L12C1B3-125-TN
% Mức: TH
% ID: L12C1B3-125-TN
% ID: L12C1B3-125-TN
% Mức: TH
% ID: L12C1B1-02-TN
%=== Câu 5 ===%
\begin{ex}
% ID: L12C1B1-02-TN
% Mức: TH
	Một lượng nước đá đang ở $0^\circ C$. Trong suốt quá trình nóng chảy của nước đá thì nhiệt độ của nó
	\choice 
	{tăng lên sau đó giảm xuống}
	{\True không thay đổi}
	{luôn giảm}
	{luôn tăng}
	\loigiai{
		Nước đá đang ở nhiệt độ nóng chảy $0^\circ C$. Trong suốt quá trình nóng chảy, nhiệt lượng cung cấp dùng để làm nước đá chuyển từ thể rắn sang thể lỏng, không dùng để tăng nhiệt độ.
		
		Do đó, \textbf{nhiệt độ của nước đá trong suốt quá trình nóng chảy không thay đổi}, vẫn bằng $0^\circ C$.
	}
\end{ex}
\dangbt{Quá trình chuyển pha}
\begin{ex}
\nguon{https://tuyensinh247.com/bai-tap-783241.html}
Vào mùa đông, khi hà hơi vào một gương phẳng, mặt gương bị mờ đi. Sự mờ đi của mặt gương được giải thích bằng hiện tượng
\choice
{nóng chảy}
{bay hơi}
{ngưng kết}
{\True ngưng tụ}
\loigiai{
Hơi nước có trong hơi thở khi gặp bề mặt gương lạnh sẽ bị ngưng tụ thành các hạt nước li ti bám trên mặt gương, làm gương bị mờ đi.
}
\end{ex}
\dangbt{Quá trình chuyển pha}
\begin{ex}
Thể nào dưới đây không thuộc ba thể cơ bản của chất?
\choice
{Thể rắn}
{Thể lỏng}
{Thể khí}
{\True Thể dẻo}
\nguon{https://hoc24.vn/trac-nghiem/bai-10-cac-the-cua-chat-va-su-chuyen-the.17035}
\loigiai{
Ba thể cơ bản của chất thường gặp là thể rắn, thể lỏng và thể khí. Thể dẻo không thuộc ba thể cơ bản này.
}
\end{ex}

\dangbt{Quá trình chuyển pha}
\begin{ex}
Sự chuyển thể nào sau đây diễn ra tại một nhiệt độ xác định?
\choice
{Ngưng tụ}
{Hoá hơi}
{\True Sôi}
{Bay hơi}
\nguon{https://hoc24.vn/trac-nghiem/bai-10-cac-the-cua-chat-va-su-chuyen-the.17035}
\loigiai{
Sự sôi là quá trình hoá hơi xảy ra cả ở bên trong và trên bề mặt khối chất lỏng tại một nhiệt độ xác định (nhiệt độ sôi). Sự bay hơi nói chung diễn ra ở nhiệt độ bất kì.
}
\end{ex}

\dangbt{Quá trình chuyển pha}
\begin{ex}
Một số chất khí có mùi thơm toả ra từ bông hoa hồng làm ta có thể ngửi thấy mùi hoa thơm. Điều này thể hiện tính chất nào của thể khí?
\choice
{Dễ dàng nén được}
{Không có hình dạng xác định}
{\True Có thể lan toả trong không gian theo mọi hướng}
{Không chảy được}
\nguon{https://hoc24.vn/trac-nghiem/bai-10-cac-the-cua-chat-va-su-chuyen-the.17035}
\loigiai{
Các phân tử khí gây mùi hương chuyển động hỗn loạn và lan toả tự do theo mọi hướng trong không gian đến khứu giác.
}
\end{ex}

\dangbt{Quá trình chuyển pha}
\begin{ex}
Hiện tượng tự nhiên nào sau đây là do hơi nước ngưng tụ?
\choice
{\True Tạo thành mây}
{Gió thổi}
{Mưa rơi}
{Lốc xoáy}
\nguon{https://hoc24.vn/trac-nghiem/bai-10-cac-the-cua-chat-va-su-chuyen-the.17035}
\loigiai{
Hơi nước bốc lên cao gặp lạnh ngưng tụ lại thành các giọt nước li ti tạo thành mây.
}
\end{ex}

\dangbt{Quá trình chuyển pha}
\begin{ex}Xét các phát biểu sau về đặc điểm các thể của chất:
\choiceTF
{\True Không khí chiếm đầy khoảng không gian xung quanh ta là vì chất khí lan truyền trong không gian theo mọi hướng.}
{\True Ta có thể bơm không khí vào trong lốp xe cho tới khi lốp xe căng lên vì chất khí có thể nén được.}
{\True Ta có thể rót nước lỏng vào cốc vì chất lỏng có thể rót được và chảy tràn trên bề mặt.}
{\True Gõ nhẹ thước kẻ vào mặt bàn thấy cả hai đều không biến dạng vì chất rắn có hình dạng cố định.}
\nguon{https://hoc24.vn/trac-nghiem/bai-10-cac-the-cua-chat-va-su-chuyen-the.17035}
\loigiai{
a đúng vì chất khí lan truyền tự do theo mọi hướng;
 b đúng vì chất khí có khoảng cách phân tử lớn nên dễ nén;
c đúng vì chất lỏng có tính linh động, chảy tràn được;
d đúng vì vật rắn có hình dạng xác định khó bị biến dạng khi tác dụng lực nhẹ.
}
\end{ex}

\dangbt{Quá trình chuyển pha}
\begin{ex}Khi nói về sự chuyển thể của các chất, xét các phát biểu sau:
\choiceTF
{\True Quá trình chất chuyển từ thể rắn sang thể lỏng gọi là sự nóng chảy.}
{\True Quá trình chất chuyển từ thể lỏng sang thể rắn gọi là sự đông đặc.}
{\True Quá trình nóng chảy của một chất tinh khiết diễn ra ở nhiệt độ xác định gọi là nhiệt độ nóng chảy.}
{\True Với cùng một chất tinh khiết, nhiệt độ nóng chảy và nhiệt độ đông đặc xảy ra ở cùng một nhiệt độ.}
\nguon{https://hoc24.vn/trac-nghiem/bai-10-cac-the-cua-chat-va-su-chuyen-the.17035}
\loigiai{
a đúng theo định nghĩa sự nóng chảy;
 b đúng theo định nghĩa sự đông đặc;
 c đúng vì mỗi chất tinh khiết nóng chảy ở một nhiệt độ xác định;
 d đúng vì với mỗi chất, nhiệt độ nóng chảy bằng nhiệt độ đông đặc.
}
\end{ex}
\dangbt{Quá trình chuyển pha}
\begin{ex}
Khi đun nóng một lượng nước đá, đồ thị nhiệt độ $t$ phụ thuộc vào thời gian đun $T$ được biểu diễn như hình vẽ.
\begin{center}
\begin{tikzpicture}[>=stealth,scale=0.9]
	% Trục
	\draw[->] (0,0) -- (8,0) node[below] {$T$ (thời gian)};
	\draw[->] (0,0) -- (0,5.5) node[left] {$t\,({}^\circ\mathrm{C})$};

	% Mức nhiệt độ
	\draw[dashed,gray] (0,1.2) -- (7.6,1.2);
	\draw[dashed,gray] (0,4.2) -- (7.6,4.2);

	\node[left] at (0,1.2) {$0^\circ\mathrm{C}$};
	\node[left] at (0,4.2) {$100^\circ\mathrm{C}$};

	% Đường biểu diễn
	\draw[very thick]
	(0.8,0.7) -- (2.4,1.2)
	-- (4.2,4.2)
	-- (6.0,4.2)
	-- (7.4,5.1);

	% Điểm A B C D E
	\foreach \x/\y in {
		0.8/0.7,
		2.4/1.2,
		4.2/4.2,
		6.0/4.2,
		7.4/5.1}
		\fill (\x,\y) circle (2pt);

	\node[below left] at (0.8,0.7) {$A$};
	\node[below] at (2.4,1.2) {$B$};
	\node[above left] at (4.2,4.2) {$C$};
	\node[above] at (6.0,4.2) {$D$};
	\node[above right] at (7.4,5.1) {$E$};
\end{tikzpicture}
\end{center}
Trong các phát biểu sau, phát biểu nào là đúng?
\choice
{Đoạn $AB$ ứng với quá trình nước đá đang tan.}
{Đoạn $BC$ ứng với quá trình nước đang sôi.}
{\True Đoạn $CD$ ứng với quá trình nước đang sôi.}
{Đoạn $DE$ ứng với quá trình nước đang bay hơi.}
\loigiai{Đoạn $AB$ ứng với quá trình nước đá đang nóng chảy. Đoạn $BC$ ứng với quá trình nước đang nóng lên. Đoạn $CD$ ứng với quá trình nước đang sôi. Đoạn $DE$ ứng với quá trình hơi nước đang nóng lên.}
\nguon{https://www.scribd.com/document/1058018453/S\%E1\%BB\%B1-Chuy\%E1\%BB\%83n-Th\%E1\%BB\%83-Bai-T\%E1\%BA\%ADp}
\end{ex}
